% SHARD-ID: AQM-71-QSA-SYMPLECTIC-FIELD-LABELS-SET
% SHARD-TITLE: Finite Symplectic Field Labels on a Set
% SHARD-SUMMARY: Reviews the finite symplectic field-label association from AQM-14 in QSA Step 07 language.
% SHARD-SUMMARY: Records the Map(X,V), finite-support, pointwise-pairing, radical-quotient, and Weyl-carrier ingredients with local anchors.
% SHARD-SUMMARY: Lists only checked \(\mathbb F_3\) CSV examples and states that new examples require new run or source rows.
% SHARD-KEYWORDS: quantum-system association, finite set, symplectic labels, Map(X,V), radical quotient, Weyl-Heisenberg, qudits

\section{Finite Symplectic Field Labels on a Set}
\label{sec:qsa-symplectic-field-labels-set}

\subsection{Status and Local Anchors}

This is QSA Step 07 from
\path{docs/quantum_system_associations/PLAN.md}.  It is a review shard, not a
new construction shard.  It adds no new examples, no new source claims, and no
new run artifacts.

The vocabulary and evidence discipline are those of
\texttt{CONVENTIONS.md} convention \texttt{(ap)} and AQM-65.  The finite
field-label convention is \texttt{CONVENTIONS.md} convention \texttt{(q)}.
The main report anchors are AQM-14, especially its dependencies D1--D4,
lemmas L1--L5, and Weyl step L6; the exact
\(\mathbb F_3\) run bundle
\path{runs/2026-05-24-arithmetic-quantum-fields/}; and AQM-22, especially
its all-map prime-field construction D3--D6 and theorems T1--T2.  The
finite-support variant is the closed-point finite-residue version recorded in
AQM-28 D5--T1 and \texttt{CONVENTIONS.md} convention \texttt{(y)}.

\subsection{Association Data}

The finite-set version starts with a finite point set \(X\), a finite
symplectic vector space \((V,\omega_V)\) over \(\mathbb F_q\), and a chosen
\(\mathbb F_q\)-linear subspace
\[
  E\leq \operatorname{Map}(X,V).
\]
As recorded in convention \texttt{(q)} and AQM-14 D1--D4, the neutral
function-space notation for a set is \(\operatorname{Map}(X,V)\).  The word
\(\operatorname{Hom}\) is reserved for a separately specified
structure-preserving class of maps.

Thus this association is not attached to a bare set alone.  The set supplies
labels for points; the workflow also needs the target \(V\), the subspace
\(E\), and the pointwise pairing weights.  In the existing checked run all
weights are \(1\).

For infinite closed-point layers over affine schemes, the finite-qudit variant
uses finite support rather than all maps.  AQM-28 defines, for an open
\(U\subset X\),
\[
  E(U)=\bigoplus_{x\in U\cap |X|_{\rm cl}}\kappa(x)^2
\]
as a finite-support direct sum over closed points.  Regular-function profiles
may have infinite support there, so AQM-28 keeps them separate from
quasi-local Weyl operators.

\subsection{Pointwise Pairing and Radical}

For \(E\leq\operatorname{Map}(X,V)\), AQM-14 and convention \texttt{(q)} use
the pointwise alternating form
\[
  \Omega_E(\phi,\psi)
  =
  \sum_{x\in X} w_x\,\omega_V(\phi(x),\psi(x)).
\]
AQM-14 L1--L2 check the pointwise vector-space and bilinear alternating-form
properties.  AQM-14 L3 records nondegeneracy for the all-map case when
\(\omega_V\) is nondegenerate and all weights are nonzero.

For a chosen subspace \(E\), the radical must still be checked:
\[
  \operatorname{rad}(E)
  =
  \{\phi\in E:\Omega_E(\phi,\psi)=0\text{ for all }\psi\in E\}.
\]
AQM-14 L4--L5 and convention \texttt{(q)} fix the rule:
\[
  E_{\rm red}=E/\operatorname{rad}(E).
\]
If the radical is zero, \(E\) itself is the Weyl label space.  If the radical
is nonzero, the honest finite symplectic Weyl label space is the quotient
\(E_{\rm red}\), not the unreduced \(E\).

In the checked \(\mathbb F_3\) scalar-function examples, AQM-14 uses
\(V=\mathbb F_3^2\) with
\[
  \omega((q,p),(q',p'))=p q'-q p',
\]
and a scalar subspace \(U\leq\operatorname{Map}(X,\mathbb F_3)\) gives
\[
  E=Uq\oplus Up,\qquad
  \Omega((q,p),(q',p'))=B(p,q')-B(q,p'),
\]
where \(B(f,g)=\sum_x f(x)g(x)\).

\subsection{Weyl Labels and Qudit Carriers}

After reduction, AQM-14 L6 chooses a symplectic basis
\[
  E_{\rm red}\cong \mathbb F_q^n\oplus\mathbb F_q^n
\]
and represents the labels on
\[
  \mathcal H=\mathbb C^{\mathbb F_q^n}\cong \ell^2(\mathbb F_q^n)
\]
by shifts, phases, and Weyl operators
\[
  T_q|a\rangle=|a+q\rangle,\qquad
  R_p|a\rangle=\chi(p\cdot a)|a\rangle,\qquad
  W(q,p)=T_qR_p.
\]
The resulting observable algebra is the complex span of these Weyl operators,
as stated in AQM-14 after L6.

AQM-22 is the prime-field qudit comparison.  It separates the one-point
spectrum \(X_{\rm sp}=\Spec\mathbb F_p\) from the \(p\)-element set
\(X_{\rm el}=\mathbb F_p\).  For \(X_{\rm el}\) and
\[
  E_{\rm el}=\operatorname{Map}(X_{\rm el},\mathbb F_p^2),
\]
AQM-22 T1 identifies the generated Weyl operators with the ambient
odd-\(p\) \(p\)-qudit Pauli group, one qudit for each element of
\(\mathbb F_p\).  AQM-22 T2 records that stabilizer codes are extra data:
one must choose an isotropic label subspace and compatible phases.

\subsection{Checked \texorpdfstring{\(\mathbb F_3\)}{F3} Rows}

The exact run
\path{runs/2026-05-24-arithmetic-quantum-fields/} was produced by
\path{scripts/bridges/arithmetic_quantum_fields.jl}.  Its README states that
the run checks pointwise symplectic function spaces, radicals, reduced Weyl
labels, Hilbert dimensions, and observable-algebra dimensions.  Its schema is
the \texttt{arithmetic\_quantum\_field\_examples.csv} entry in
\path{data/SCHEMA.md}.

The table below is only a review of existing CSV rows in
\path{runs/2026-05-24-arithmetic-quantum-fields/data/arithmetic_quantum_field_examples.csv}.
It is not a new computation.

\begin{center}
\small
\begin{tabular}{lrrrrr}
\toprule
CSV row & \(|X|\) & \(\dim U\) & \(\dim\operatorname{rad}\) & \(|E_{\rm red}|\) & \(\dim\mathcal H\)\\
\midrule
\texttt{finite\_two\_points\_all\_maps} & 2 & 2 & 0 & 81 & 9\\
\texttt{finite\_set\_3\_full} & 3 & 3 & 0 & 729 & 27\\
\texttt{vector\_space\_F3\_affine} & 3 & 2 & 2 & 9 & 3\\
\texttt{affine\_plane\_F3\_linear} & 9 & 2 & 4 & 1 & 1\\
\texttt{parabola\_F3\_degree\_le\_1} & 3 & 3 & 0 & 729 & 27\\
\bottomrule
\end{tabular}
\end{center}

The companion basis file
\path{runs/2026-05-24-arithmetic-quantum-fields/data/arithmetic_quantum_field_bases.csv}
records the scalar basis values and scalar Gram rows behind those summary
rows.  For example, the row group
\texttt{vector\_space\_F3\_affine} records basis labels \(1,t\) and Gram rows
\texttt{0 0} and \texttt{0 2}; the row group
\texttt{parabola\_F3\_degree\_le\_1} records basis labels \(1,x,y\) and Gram
rows \texttt{0 0 2}, \texttt{0 2 0}, and \texttt{2 0 2}.

\subsection{Boundary}

This shard closes only the review requested by QSA Step 07.  It does not
extend the example list beyond the existing AQM-14/AQM-22 material and the
2026-05-24 CSV rows.  Any new finite set, field space, weighting, radical
calculation, qudit carrier count, or comparison with another association
workflow would require a new local source, derivation, or run row before it
can be stated as report evidence.
